\documentclass[aspectratio=169,11pt]{beamer}
\usepackage[T1]{fontenc}
\usepackage{lmodern,amsmath,booktabs,graphicx}
\definecolor{ink}{HTML}{183248}
\definecolor{teal}{HTML}{187B84}
\definecolor{rust}{HTML}{CB593C}
\setbeamercolor{normal text}{fg=ink,bg=white}
\setbeamercolor{frametitle}{fg=ink}
\setbeamercolor{structure}{fg=teal}
\setbeamercolor{alerted text}{fg=rust}
\setbeamerfont{frametitle}{size=\Large,series=\bfseries}
\setbeamertemplate{navigation symbols}{}
\setbeamertemplate{footline}{\hfill\scriptsize Siyan Wang\enspace|\enspace Preliminary reproduction\enspace|\enspace\insertframenumber/\inserttotalframenumber\hspace{5mm}\vspace{3mm}}
\setbeamersize{text margin left=7mm,text margin right=7mm}
\newcommand{\src}[1]{\par\vspace{1mm}{\fontsize{6.5}{8}\selectfont\color{gray}#1\par}}
\title{9/30 research meeting}
\author{Siyan Wang}
\date{30 September 2026}

\begin{document}

% ==============================================================================
% Title page — plain frames show no footer and do not increment slide numbers.
% ==============================================================================
\begin{frame}[plain,noframenumbering]
  \titlepage
\end{frame}

\section{European heatwave attribution: reproduction results}

\begin{frame}[plain,noframenumbering]
  \sectionpage
\end{frame}


% ==============================================================================
% Slide 01 — Aim and workflow
% Purpose: state the reproduction question, summarize the full pipeline, and interpret long return periods.
% ================================================================================
\begin{frame}{Aim and reproduction workflow}
\small
\textbf{Aim}\par\smallskip
\footnotesize Reproduce the report's observation-based attribution of European heat extremes, test robustness across ERA5, CPC, E-OBS and Berkeley Earth, and compare our fitted curves and attribution estimates with the published Figure~7 and Table~2.
\medskip

\textbf{Workflow}\par\smallskip
\scriptsize
\begin{beamercolorbox}[sep=1.5mm,rounded=true]{block body}
Daily gridded Tmax/Tmin
$\rightarrow$ fixed land footprint
$\rightarrow$ area-weighted regional daily temperature
$\rightarrow$ three-day mean
$\rightarrow$ Annual/June Tx3x or Tn3x
$\rightarrow$ pair with smoothed NASA GMST
$\rightarrow$ fit a nonstationary GEV
$\rightarrow$ compare factual and cooler climates.
\end{beamercolorbox}
\medskip

\textbf{How can a record shorter than 100 years give a 100-year return period?}\par\smallskip
\footnotesize All available annual extremes estimate one GEV distribution. At the 2026 factual GMST, the 100-year return level is its 99th percentile:
\[
P(X>z_{100}\mid G_{2026})=0.01,
\qquad
z_{100}=F^{-1}_{\mathrm{GEV}}(0.99\mid G_{2026}).
\]
A ``100-year event'' therefore means a fitted annual exceedance probability of 1\%; it does not require observing a complete 100-year cycle. When the record is shorter than 100 years, this is tail extrapolation, so it is sensitive to the GEV shape and has wider uncertainty; we report bootstrap intervals alongside the point estimate.
\src{Workflow follows the report's observational attribution design and the reproduction scripts.}
\end{frame}


% ================================================================================
% Slide 02 — Data products
% Purpose: distinguish reanalysis from observation-based gridded products.
% ==============================================================================
\begin{frame}{Four gridded temperature products}
\footnotesize
\renewcommand{\arraystretch}{1.12}
\begin{tabular}{p{1.25cm}p{4.35cm}p{2.15cm}p{5.35cm}}\toprule
Product & Type of data & Fit period & Main characteristic\\\midrule
ERA5 & Reanalysis: observations assimilated into a numerical weather model & 1950--2025 & Continuous and physically consistent; model-dependent.\\
E-OBS & Gridded station observations over European land (spatial interpolation) & 1920--2025 & High European resolution with uncertainty estimates; coverage varies by country and period.\\
CPC & Gridded station-based daily temperature analysis (GTS observations) & 1979--2025 & Frequently updated, but has the shortest record.\\
Berkeley & Observation-based gridded daily land temperature reconstruction & 1897--2023 annual; 1897--2024 June & Longest record with breakpoint correction; daily data are experimental.\\\bottomrule
\end{tabular}
\medskip
\scriptsize All four products are gridded. ERA5 combines observations with a numerical model. The other three construct spatial fields primarily from station observations. Agreement across them tests sensitivity to different source data and processing methods.
\src{Sources: report Appendix A1.1; product metadata; saved fit-year records.}
\end{frame}


% ==============================================================================
% Slide 03 — NASA GISTEMP upstream construction
% Purpose: explain where one global monthly anomaly comes from before our code starts.
% ================================================================================
\begin{frame}{NASA turns many observations into one global anomaly each month}
\footnotesize
\begin{columns}[T,onlytextwidth]
\column{.48\textwidth}
\textbf{1. Two monthly input fields}
\begin{itemize}
  \item \textbf{Land:} adjusted monthly mean surface-air temperatures from NOAA GHCN v4 stations.
  \item \textbf{Ocean:} monthly sea-surface temperatures from NOAA ERSST v5.
\end{itemize}
\medskip
\textbf{2. Work with local anomalies}
\begin{itemize}
  \item For each location and calendar month, subtract its own 1951--1980 mean.
  \item January is compared with January, February with February, and so on; this removes local climate and the seasonal cycle.
\end{itemize}

\column{.48\textwidth}
\textbf{3. Build a spatially complete index}
\begin{itemize}
  \item GISTEMP processes the anomaly series on an equal-area grid; its standard land analysis uses 1200-km smoothing to estimate nearby unsampled areas.
  \item Land anomalies are combined with the ERSST ocean anomaly field.
\end{itemize}
\medskip
\textbf{4. Average globally}
\begin{itemize}
  \item Equal-area subboxes are aggregated into latitude zones and broad latitude bands.
  \item Area-weighting those bands gives one Land--Ocean Temperature Index (L-OTI) anomaly for that month.
\end{itemize}
\end{columns}
\begin{beamercolorbox}[sep=1.5mm,rounded=true]{block body}
\scriptsize \textbf{Starting point of this reproduction:} NASA has already performed the station/ocean merging, gap handling and global spatial averaging. We download the resulting monthly global index; we do not average raw stations ourselves.
\end{beamercolorbox}
\src{Sources: NASA GISTEMP v4 data page and FAQ; GHCN v4 + ERSST v5; 1951--1980 anomaly baseline.}
\end{frame}


% ================================================================================
% Slide 04 — Downloaded GMST structure and our annualization
% Purpose: separate NASA's upstream global aggregation from our temporal processing.
% ================================================================================
\begin{frame}{The NASA output is a year-by-month table; our code annualizes it}
\begin{columns}[T,onlytextwidth]
\column{.49\textwidth}
\textbf{What we download}\par\smallskip
\scriptsize
\setlength{\tabcolsep}{3.2pt}
\begin{tabular}{rrrrrrr}\toprule
Year & Jan & Feb & Mar & $\cdots$ & Dec & J--D\\\midrule
2023 & 0.88 & 0.97 & 1.23 & $\cdots$ & 1.36 & 1.17\\
2024 & 1.25 & 1.44 & 1.40 & $\cdots$ & 1.27 & 1.29\\
2025 & 1.38 & 1.26 & 1.37 & $\cdots$ & 1.06 & 1.19\\
2026 & 1.09 & 1.25 & 1.32 & $\cdots$ & missing & missing\\\bottomrule
\end{tabular}
\par\medskip
\footnotesize One row is a year; Jan--Dec are NASA global monthly L-OTI anomalies in $^\circ$C. ``J--D'' is NASA's annual summary, but our code reads the twelve monthly columns directly.

\column{.47\textwidth}
\textbf{Our Step 1: one value per complete year}\par
\small
\[
 g_y=\frac{1}{12}\sum_{m=1}^{12}g_{y,m}.
\]
\footnotesize We calculate $g_y$ only when all twelve NASA monthly global anomalies are finite. Thus a partial 2026 row is not treated as an observed annual GMST.
\medskip

\textbf{Our Step 2: the report's four-year smoothing}\par
\small
\[
G_t=\frac14(g_{t-2}+g_{t-1}+g_t+g_{t+1}).
\]
\footnotesize The four-year mean is assigned to year 3 of its window, so $G_{2026}$ requires annual values for 2024--2027.
\end{columns}
\src{Sources: NASA GISTEMP v4 global L-OTI CSV; prepare\_gmst.R.}
\end{frame}


% ================================================================================
% Slide 05 — Smoothed GMST endpoint
% Purpose: make the provisional endpoint calculation and its status explicit.
% ================================================================================
\begin{frame}{The 2026 smoothed endpoint requires two provisional annual values}
\begin{columns}[T,onlytextwidth]
\column{.50\textwidth}
\textbf{Why extrapolation enters}\par\medskip
\small NASA's 2026 row is incomplete and 2027 has not been observed, but the year-3 four-year window needs both annual values. For this reproduction, we fit
\[
 g_y=a+by+\varepsilon_y
\]
by ordinary least squares to the last fifteen complete annual anomalies, 2011--2025, and predict 2026 and 2027.

\column{.46\textwidth}
\textbf{Values used at the endpoint}\par\smallskip
\footnotesize
\begin{tabular}{rrrr}\toprule
2024 & 2025 & 2026* & 2027*\\\midrule
1.28583 & 1.19333 & 1.22714 & 1.26606\\\bottomrule
\end{tabular}
\[
\widehat G_{2026}
=1.24309^\circ\mathrm C.
\]
\scriptsize *OLS predictions, not observations.
\end{columns}
\begin{beamercolorbox}[sep=1.5mm,rounded=true]{alertblock body}
\scriptsize The report specifies the four-year, year-3 smoothing convention but does not document this endpoint extrapolation. Therefore $\widehat G_{2026}$ is a provisional reproduction assumption.
\end{beamercolorbox}
\src{Sources: report Appendix A1.1; prepare\_gmst.R; saved endpoint assumptions.}
\end{frame}


% ================================================================================
% Slide 06 — Why anomalies work
% Purpose: answer why anomaly data can act as the GMST covariate.
% ==============================================================================
\begin{frame}{A fixed anomaly baseline does not change the attribution fit}
\begin{columns}[T,onlytextwidth]
\column{.52\textwidth}
\small Suppose absolute GMST and the reported anomaly differ by a constant $c$:
\[
G_t^{\mathrm{absolute}}=G_t^{\mathrm{anomaly}}+c.
\]
The report's location model can be written as
\[
\mu_t=\mu_0+\alpha G_t^{\mathrm{anomaly}}
     =(\mu_0-\alpha c)+\alpha G_t^{\mathrm{absolute}}.
\]
Changing the baseline changes only the intercept.

\column{.44\textwidth}
\textbf{Quantities that remain unchanged}\par\medskip
\small
\begin{itemize}
  \item the GMST slope $\alpha$;
  \item differences such as $G_F-G_C$;
  \item fitted probabilities and return levels;
  \item probability ratios and intensity changes.
\end{itemize}
\medskip
\footnotesize Our code centers GMST internally at its fit-period mean for numerical convenience. That is an algebraically equivalent parameterization, not a separate method specified by the report.
\end{columns}
\src{Source: report Appendix A1.2; helpers/gev.R.}
\end{frame}


% ==============================================================================
% Slide 07 — Origins and upstream processing of temperature products
% Purpose: distinguish upstream product construction from our dataset-specific choices.
% ================================================================================
\begin{frame}{How the four daily temperature products are constructed}
\scriptsize
\setlength{\tabcolsep}{3pt}
\renewcommand{\arraystretch}{1.20}
\begin{tabular}{p{1.05cm}p{2.75cm}p{4.35cm}p{5.35cm}}\toprule
Product & Input observations & Upstream product construction & Choice made in this reproduction\\\midrule
ERA5 & Multi-source observations + numerical weather prediction
& Data assimilation creates a physically consistent analysis; short forecasts supply hourly maximum 2-m temperature (`mx2t').
& Climate Explorer derives the daily maximum; we convert K to $^\circ$C.\\
CPC & GTS station observations
& Station quality control followed by interpolation to a regular global grid.
& We download the regional daily subset; values are already in $^\circ$C.\\
E-OBS & European station observations
& Interpolation over European land; an ensemble describes interpolation uncertainty.
& We use the ensemble mean from v33.0e; the report used v32.0e.\\
Berkeley & Merged station observations
& Breakpoint correction followed by spatial reconstruction of the anomaly field.
& We use `full': anomaly plus daily climatology, yielding absolute Tmax.\\\bottomrule
\end{tabular}
\medskip
\footnotesize The products therefore differ before our code begins: ERA5 is model--observation reanalysis, while the other three reconstruct gridded fields mainly from stations. Our later workflow applies a common regional aggregation while retaining these upstream differences.
\src{Sources: NetCDF metadata and histories; product documentation; report Appendix A1.1.}
\end{frame}


% ================================================================================
% Slide 08 — Raw file structures
% Purpose: show enough of the actual arrays to make later aggregation concrete.
% ==============================================================================
\begin{frame}{The raw temperature files are daily space--time arrays}
\footnotesize
\renewcommand{\arraystretch}{1.15}
\begin{tabular}{p{1.45cm}p{3.0cm}p{2.1cm}p{2.0cm}p{2.15cm}p{2.4cm}}\toprule
Product & Variable & Stored dimensions & Regional grid & Unit & Year starts\\\midrule
ERA5 & `tmax(time,lat,lon)' & $28002\times81\times123$ & about $0.25^\circ$ & K & 1950\\
CPC Tmax/Tmin & `tmax(time,lat,lon)' & $17423\times40\times61$ & $0.5^\circ$ & $^\circ$C & 1979\\
E-OBS & `tx'/`tn(time,lat,lon)' & $38929\times80\times122$ & $0.25^\circ$ & $^\circ$C & 1979\\
Berkeley & `TMAX(time,lat,lon)' & $46617\times20\times31$ & $1^\circ$ & $^\circ$C & 1897 \\\bottomrule
\end{tabular}
\medskip
\small Each day is a two-dimensional temperature map over $40$--$60^\circ$N, $10.5^\circ$W--$20^\circ$E. The task is to convert each map into one land-only regional temperature without allowing the represented area to change arbitrarily over time.
\src{Source: headers of the six downloaded NetCDF files. Dimensions are shown as time $\times$ latitude $\times$ longitude.}
\end{frame}


% ==============================================================================
% Slide 09 — Why a footprint enters the workflow
% Purpose: connect the report's event box to the need for a fixed land-support mask.
% ================================================================================
\begin{frame}{The event box does not yet define what enters the regional mean}
\small
\textbf{The report first defines the event in space:}
$10.5^\circ$W--$20^\circ$E and $40^\circ$--$60^\circ$N.
That rectangle is the common outer boundary for all four products.
\medskip

\begin{columns}[T,onlytextwidth]
\column{.30\textwidth}
\begin{block}{Inside the box}
\footnotesize There are land cells, ocean cells and coastal cells. The four products do not mark or observe them in the same way.
\end{block}

\column{.30\textwidth}
\begin{block}{If we average blindly}
\footnotesize ERA5 would mix land and ocean, while station-based products would average only whichever cells happen to be available on each day.
\end{block}

\column{.32\textwidth}
\begin{block}{Required before averaging}
\footnotesize We must decide once which cells represent the study region, then hold that spatial support fixed through time.
\end{block}
\end{columns}
\vspace{1mm}
\[
\boxed{\text{event bounding box}\;\longrightarrow\;\text{fixed eligible cells}\;\longrightarrow\;\text{daily regional mean}}
\]
\begin{beamercolorbox}[sep=1.5mm,rounded=true]{block body}
\footnotesize \textbf{Footprint} means this fixed set of eligible grid cells. It is a data-processing mask, not the observed outline of the 2026 heatwave and not the set of cells that were hottest on a particular day.
\end{beamercolorbox}
\src{Source: report Section 1.3 for the event box; reproduction workflow for the fixed spatial support.}
\end{frame}


% ================================================================================
% Slide 10 — Dataset-specific footprint construction
% Purpose: show how the fixed support is built and why ERA5 differs.
% ================================================================================
\begin{frame}{We construct the footprint once, then reuse it every day}
\begin{columns}[T,onlytextwidth]
\column{.49\textwidth}
\textbf{CPC, E-OBS and Berkeley}\par\medskip
\small These station-based grids generally lack values over ocean. We use persistent data availability as a practical proxy for stable land coverage:
\[
M_i=\mathbf 1\!\left(
\frac{\#\{d:T_{d,i}\ \mathrm{finite}\}}{\#\{d\}}\ge0.95
\right).
\]
\footnotesize A cell is eligible only when it contains a finite value on at least 95\% of all dates in that product's downloaded record. The resulting $M_i$ is fixed for the full daily calculation.

\column{.47\textwidth}
\textbf{ERA5}\par\medskip
\small ERA5 supplies valid 2-m temperature over both land and ocean. Its own missingness therefore cannot identify land. We map the saved E-OBS footprint to the nearest ERA5 grid cells and use that mapped mask on every ERA5 day.
\medskip

\textbf{Interpretation}\par\smallskip
\footnotesize This makes the four regional series approximately land-comparable and prevents daily spatial coverage from drifting. It is a reproduction choice: the 95\% rule mixes coastline with long-record availability, and nearest-neighbour mapping can misclassify coastal ERA5 cells. A native ERA5 land--sea mask is a useful sensitivity test.
\end{columns}
\src{Sources: helpers/workflow.R; Code/run\_era5\_tx.R; saved grid\_footprint.csv files.}
\end{frame}


% ================================================================================
% Slide 11 — Spatial aggregation
% Purpose: show exactly how a daily grid becomes one regional value.
% ==============================================================================
\begin{frame}{Each daily grid becomes one area-weighted regional value}
\begin{columns}[T,onlytextwidth]
\column{.54\textwidth}
\small For day $d$, grid cell $i$ has latitude $\phi_i$, mask $M_i$, and validity indicator $V_{d,i}$. We use
\[
w_i=\cos(\phi_i)M_i,
\]
\[
\overline T_d=
\frac{\sum_i V_{d,i}w_iT_{d,i}}
     {\sum_i V_{d,i}w_i}.
\]
The cosine weight approximates the area represented by a regular longitude--latitude cell. The denominator is recalculated from valid cells each day.

\column{.42\textwidth}
\textbf{Daily quality rule}\par\medskip
\small Define weighted coverage as
\[
C_d=\frac{\sum_iV_{d,i}w_i}{\sum_iw_i}.
\]
If $C_d<0.95$, we set the regional daily value to missing.\par\medskip
\footnotesize This step produces one regional daily Tmax or Tmin. It does not select the hottest grid cell and does not average land and ocean together.
\end{columns}
\src{Source: helpers/workflow.R and the dataset-specific run scripts.}
\end{frame}


% ==============================================================================
% Slide 12 — Three-day extremes
% Purpose: explain order of operations and the meaning of Tx3x/Tn3x.
% ==============================================================================
\begin{frame}{Regional daily temperature becomes Tx3x or Tn3x}
\begin{columns}[T,onlytextwidth]
\column{.55\textwidth}
\textbf{Common processing order}\par\medskip
\small
\begin{enumerate}
  \item Calculate the land-only regional value $\overline T_d$ for each day.
  \item Calculate a trailing three-day mean:
  \[
  \overline T_{3,d}=\frac{\overline T_{d-2}+\overline T_{d-1}+\overline T_d}{3}.
  \]
  \item Take the maximum within each calendar year or within June.
\end{enumerate}

\column{.41\textwidth}
\textbf{Definitions}\par\medskip
\small
\[
\mathrm{Tx3x}_y=\max_d\overline T^{\max}_{3,d},
\]
\[
\mathrm{Tn3x}_y=\max_d\overline T^{\min}_{3,d}.
\]
\footnotesize Tn3x is the maximum three-day mean of daily minimum temperature, representing the warmest sequence of nights. It is not a minimum.\par\medskip
An Annual or June block enters the fit only if at least 95\% of its daily regional values are finite. June windows remain fully inside June; annual windows remain inside the calendar year.
\end{columns}
\src{Source: helpers/workflow.R. We first average in space, then average over three days, then take the temporal maximum.}
\end{frame}


% ==============================================================================
% Slide 13 — GEV and finite-record return levels
% Purpose: explain how a 100-year threshold can be estimated from shorter data.
% ==============================================================================
\begin{frame}{A 100-year threshold can exceed the length of the record}
\begin{columns}[T,onlytextwidth]
\column{.50\textwidth}
\textbf{One nonstationary GEV across all years}\par\medskip
\small Each annual or June extreme $X_t$ is paired with its smoothed GMST $G_t$:
\[
X_t\sim\mathrm{GEV}(\mu_t,\sigma,\xi),
\qquad \mu_t=\mu_0+\alpha G_t.
\]
The data estimate the common parameters $\mu_0,\alpha,\sigma,\xi$. We do not fit a separate GEV to each year.

\column{.46\textwidth}
\textbf{Return-level threshold}\par\medskip
\small For return period $T$, the factual threshold is
\[
x_T=F_F^{-1}\!\left(1-\frac1T\right).
\]
For $T=100$, this is the fitted 0.99 quantile. A 47-year CPC record can therefore yield a 100-year threshold because the parametric GEV extrapolates beyond the observed ranks.\par\medskip
\footnotesize The estimate is model-dependent and can be highly uncertain. Parameter bootstrap refits simulated records and propagates this tail uncertainty into intervals. It does not create additional observed years.
\end{columns}
\src{Source: report Appendix A1.2; helpers/gev.R and helpers/workflow.R.}
\end{frame}


% ==============================================================================
% Slide 14 — Model choice and attribution
% Purpose: define RR/intensity and record the report's ENSO/NAO decision.
% ==============================================================================
\begin{frame}{Attribution compares the same event across climate states}
\begin{columns}[T,onlytextwidth]
\column{.49\textwidth}
\textbf{Event definitions used in the report}\par\smallskip
\small
\begin{tabular}{lrr}\toprule
Return period & Tx3x & Tn3x\\\midrule
Annual & 25 yr & 20 yr\\
June & 100 yr & 100 yr\\\bottomrule
\end{tabular}

\medskip

\footnotesize
$G_F$ is the assumed 2026 GMST. Counterfactual climates use
\[
G_C=G_F-0.6^\circ\mathrm C
\quad\text{or}\quad
G_F-1.1^\circ\mathrm C,
\]
representing 2003-like and 1976-like warming levels.

\column{.47\textwidth}
\textbf{Probability and intensity changes}\par
\scriptsize
\[
x_F=F_F^{-1}(1-1/T),\qquad
RR=\frac{1/T}{1-F_C(x_F)},
\]
\[
FAR= 1-\frac{1}{RR}
\]
\[
\Delta I=F_F^{-1}(1-1/T)-F_C^{-1}(1-1/T)
=\alpha\Delta G.
\]
\footnotesize Appendix A1.2 also tests $\mu_t=\mu_0+\alpha G_t+\beta I_t$, with detrended Ni\~no3.4 or NAO as $I_t$. Table 3 finds no AIC improvement, so the report excludes these indices from the remaining analysis. We retain the GMST-only model.
\end{columns}
\src{Sources: report Appendix A1.2, Section 2.1.1 and Table 3.}
\end{frame}


% ==============================================================================
% Slide 15 — Figure 7 construction
% Purpose: explain every graphical element before showing comparisons.
% ==============================================================================
\begin{frame}{How to read Figure 7}
\begin{columns}[T,onlytextwidth]
\column{.49\textwidth}\centering
\includegraphics[width=\linewidth,height=5.6cm,keepaspectratio,trim=0 263 341 25,clip]{../materials/0930_reproduction_assets/figure7.pdf}

\column{.47\textwidth}
\scriptsize
\textbf{Curves.} For probabilities $p=1-1/T$, the red and blue curves plot $F_F^{-1}(p)$ and $F_P^{-1}(p)$ against return period $T$. Red uses 2026 GMST; blue uses $G_F-1.1^\circ$C.\par\smallskip
\textbf{Points.} Each observed annual extreme is shifted to a common climate before ranking:
\[
X_{t,F}=X_t+\widehat\alpha(G_F-G_t),\qquad
X_{t,P}=X_t+\widehat\alpha(G_P-G_t).
\]
The x-coordinate is an empirical return period from the rank, not the observation year. Red and blue points are two adjustments of the same observations.\par\smallskip
\textbf{Bands.} Shading gives 95\% intervals from parameter bootstrap fits.\par\smallskip
\textbf{Purple line and bottom marks.} The horizontal line is the current 25-year Annual Tx3x threshold. The red and blue marks show its return period in the two fitted climates.
\end{columns}
\src{Source: plot\_figure7.R. The report caption defines the climate curves and threshold; our code makes the point and band construction explicit.}
\end{frame}


% ==============================================================================
% Slide 16 — ERA5 Figure 7 comparison
% ==============================================================================
\begin{frame}{Figure 7 comparison: ERA5 annual Tx3x}
\begin{columns}[T,onlytextwidth]
\column{.47\textwidth}\centering\small\textbf{Report: Figure 7a}\par
\includegraphics[height=4.3cm,trim=104 575 303 72,clip]{/Users/wangcaiyan/extreme-event-attribution/0923_reproduction_assets/report_page12.pdf}
\column{.49\textwidth}\centering\small\textbf{Our reproduction: ERA5}\par
\includegraphics[width=\linewidth,height=4.3cm,keepaspectratio,trim=0 263 341 25,clip]{../materials/0930_reproduction_assets/figure7.pdf}
\end{columns}
\medskip\small
Both fits shift upward from the 1976-like climate to the current climate. Our fitted current 25-year level is $30.90^\circ$C; the same threshold has a fitted past return period of about $1.01\times10^6$ years.
\src{Report p.12, Fig.7a; saved ERA5 Figure 7 reproduction. Fit period: 1950--2025; plotted bands use 300 parameter-bootstrap fits.}
\end{frame}


% ==============================================================================
% Slide 17 — CPC Figure 7 comparison
% ==============================================================================
\begin{frame}{Figure 7 comparison: CPC annual Tx3x}
\begin{columns}[T,onlytextwidth]
\column{.47\textwidth}\centering\small\textbf{Report: Figure 7b}\par
\includegraphics[height=4.3cm,trim=298 575 107 72,clip]{/Users/wangcaiyan/extreme-event-attribution/0923_reproduction_assets/report_page12.pdf}
\column{.49\textwidth}\centering\small\textbf{Our reproduction: CPC}\par
\includegraphics[width=\linewidth,height=4.3cm,keepaspectratio,trim=346 263 0 25,clip]{../materials/0930_reproduction_assets/figure7.pdf}
\end{columns}
\medskip\small
Our fitted current 25-year threshold is $33.18^\circ$C. The blue curve evaluates the same fitted GEV shape and scale after shifting its location to the 1976-like GMST; the threshold then has a return period of 229 years.
\src{Report p.12, Fig.7b; saved CPC Figure 7 reproduction. The red and blue points are climate-adjusted versions of the same annual extremes.}
\end{frame}


% ==============================================================================
% Slide 18 — Berkeley Figure 7 comparison
% ==============================================================================
\begin{frame}{Figure 7 comparison: Berkeley annual Tx3x}
\begin{columns}[T,onlytextwidth]
\column{.47\textwidth}\centering\small\textbf{Report: Figure 7c}\par
\includegraphics[height=4.3cm,trim=104 390 291 264,clip]{/Users/wangcaiyan/extreme-event-attribution/0923_reproduction_assets/report_page12.pdf}
\column{.49\textwidth}\centering\small\textbf{Our reproduction: Berkeley}\par
\includegraphics[width=\linewidth,height=4.3cm,keepaspectratio,trim=0 0 341 277,clip]{../materials/0930_reproduction_assets/figure7.pdf}
\end{columns}
\medskip\small
Our current 25-year threshold is $31.92^\circ$C; its fitted 1976-like return period is 4,876 years. Berkeley has the longest record, but the far-right tail still depends on the fitted GEV shape rather than direct empirical recurrence.
\src{Report p.12, Fig.7c; saved Berkeley Figure 7 reproduction. Original axes retained; panel scales differ.}
\end{frame}


% ==============================================================================
% Slide 19 — E-OBS Figure 7 comparison
% ==============================================================================
\begin{frame}{Figure 7 comparison: E-OBS annual Tx3x}
\begin{columns}[T,onlytextwidth]
\column{.47\textwidth}\centering\small\textbf{Report: Figure 7d}\par
\includegraphics[height=4.3cm,trim=298 390 107 264,clip]{/Users/wangcaiyan/extreme-event-attribution/0923_reproduction_assets/report_page12.pdf}
\column{.49\textwidth}\centering\small\textbf{Our reproduction: E-OBS v33e}\par
\includegraphics[width=\linewidth,height=4.3cm,keepaspectratio,trim=346 0 0 277,clip]{../materials/0930_reproduction_assets/figure7.pdf}
\end{columns}
\medskip\small
Our current 25-year threshold is $32.05^\circ$C; its fitted 1976-like return period is 52,044 years. The report uses v32.0e, while our input uses v33e, so differences can reflect both sampling and product-version changes.
\src{Report p.12, Fig.7d; saved E-OBS Figure 7 reproduction. Original axes retained; panel scales differ.}
\end{frame}


% ==============================================================================
% Slide 20 — Table 2 benchmark
% Purpose: distinguish report synthesis from single-dataset estimates.
% ==============================================================================
\begin{frame}{How to read the report's Table 2}
\centering
\includegraphics[width=13.4cm,height=4.7cm,keepaspectratio,trim=44 92 59 528,clip]{/Users/wangcaiyan/extreme-event-attribution/0923_reproduction_assets/report_page12.pdf}
\par\medskip\raggedright
\small Each row fixes an event definition and return period. The two climate comparisons report a probability ratio and a change in magnitude. Parentheses are reported 95\% intervals; ``$>1$m'' means greater than $10^6$.\par\smallskip
\footnotesize Table 2 is a synthesis across observational and reanalysis evidence. Our colored estimates on the next page are individual-dataset results, so Table 2 is a benchmark rather than the expected value for each separate dataset.
\src{Original Table 2, report p.12; synthesis description on p.13.}
\end{frame}


% ==============================================================================
% Slide 21 — Intensity plot explanation
% Purpose: explain the comparison graphic and its uncertainty bars.
% ==============================================================================
\begin{frame}{Intensity comparison: estimate, interval and data source}
\centering
\includegraphics[width=14.2cm,height=5.25cm,keepaspectratio]{../materials/0930_reproduction_assets/intensity_comparison.pdf}
\par\raggedright\small
Each point estimates the warming-related intensity change from the 2003-like climate to 2026. Horizontal lines are 95\% intervals. Black diamonds reproduce the report synthesis; colored circles are our separate dataset fits. The left and right panels use Annual and June extremes.\par\smallskip
\footnotesize Overlapping intervals show that estimates are statistically compatible at this resolution; they do not prove exact reproduction. Our intervals use parameter bootstrap and omit uncertainty in GMST, the spatial mask and dataset choice.
\src{Report Table 2 and saved single-dataset estimates. Units: $^\circ$C.}
\end{frame}


% ==============================================================================
% Slide 22 — Probability-ratio comparison
% Purpose: show the tail-sensitive results and explain infinities.
% ==============================================================================
\begin{frame}{Probability-ratio comparison}
\small 
Each cell shows \textbf{2003-like / 1976-like }.

PR in the report and RR in our code denote the same probability ratio.

\medskip
\scriptsize
\setlength{\tabcolsep}{3pt}
\renewcommand{\arraystretch}{1.5}
\centering
\begin{tabular}{llrrrrr}\toprule
Index & Period & Report & ERA5 & CPC & E-OBS & Berkeley\\\midrule
Tx3x & Annual & $13.9/580$ & $37.1/40221$ & $3.74/9.15$ & $19.4/2082$ & $9.53/195$\\
Tx3x & June & $286/{>}10^6$ & $3788/\infty$ & $\infty/\infty$ & $77.3/3.09{\times}10^{17}$ & $57.3/\infty$\\
Tn3x & Annual & $163/{>}10^6$ & --- & $103/6.47{\times}10^{10}$ & $7.68/255$ & ---\\
Tn3x & June & $295/{>}10^6$ & --- & $\infty/\infty$ & $7.72/61.6$ & ---\\\bottomrule
\end{tabular}
\par\medskip\raggedright\small
Risk ratios are more sensitive than intensity changes to the fitted upper tail. $\infty$ means the fitted counterfactual exceedance probability is numerically zero, often because a negative-shape GEV has a finite upper endpoint. It does not prove physical impossibility.
\src{Sources: original Table 2 and saved PROVISIONAL result CSVs. Point estimates shown.}
\end{frame}


% ==============================================================================
% Slide 23 — Fraction of attributable risk comparison
% Purpose: express the same probability-ratio results on the bounded FAR scale.
% ================================================================================
\begin{frame}{Fraction of attributable risk (FAR) comparison}
\small
Each cell shows \textbf{2003-like / 1976-like}.

\[
\mathrm{FAR}=1-\frac{1}{\mathrm{RR}}.
\]
\vspace{-2mm}
\scriptsize
\setlength{\tabcolsep}{3pt}
\renewcommand{\arraystretch}{1.5}
\centering
\begin{tabular}{llrrrrr}\toprule
Index & Period & Report* & ERA5 & CPC & E-OBS & Berkeley\\\midrule
Tx3x & Annual & $0.9281/0.9983$ & $0.9731/0.99998$ & $0.7327/0.8907$ & $0.9484/0.99952$ & $0.8950/0.9949$\\
Tx3x & June & $0.9965/{>}0.999999$ & $0.99974/1$ & $1/1$ & $0.9871/1$ & $0.9825/1$\\
Tn3x & Annual & $0.9939/{>}0.999999$ & --- & $0.9903/1$ & $0.8698/0.9961$ & ---\\
Tn3x & June & $0.9966/{>}0.999999$ & --- & $1/1$ & $0.8704/0.9838$ & ---\\\bottomrule
\end{tabular}
\par\medskip\raggedright\small
FAR is bounded by one, so very different risk ratios can look almost identical on this scale. A displayed value of 1 results from an infinite or numerically enormous RR; it should not be interpreted as proof that the event was physically impossible in the cooler climate.
\par\smallskip
\footnotesize *The report does not tabulate FAR. Its values here are derived from the published probability ratios using $1-1/\mathrm{PR}$; ${>}10^6$ implies FAR ${>}0.999999$.
\src{Sources: original Table 2; /0930\_reproduction\_assets/results.csv. Point estimates shown.}
\end{frame}


% ================================================================================
% Slide 24 — Reproduction assessment
% Purpose: close with what is robust and what still limits exact agreement.
% ==============================================================================
\begin{frame}{Reproduction assessment and remaining gaps}
\begin{columns}[T,onlytextwidth]
\column{.48\textwidth}
\textbf{What is reproduced}\par\medskip
\small The pipeline from daily grids to regional Tx3x/Tn3x, the GMST-dependent location-shift GEV, the prescribed climate comparisons and the structure of Figure 7.\par\medskip
ERA5 Tx3x, CPC, E-OBS and Berkeley individual fits all give a positive warming signal.

\column{.48\textwidth}
\textbf{What prevents an exact match}\par\medskip
\small Product versions, land masks, missing-data rules and endpoint dates may differ from the authors' inputs.\par\medskip
Our 2026 GMST endpoint uses an undocumented extrapolation. Bootstrap intervals omit GMST, observational and mask uncertainty. Tail probability estimates remain particularly unstable.
\end{columns}
\src{Evidence: report Appendix A1; original Figure 7 and Table 2; reproduction scripts, NetCDF metadata and output CSVs.}
\end{frame}

\end{document}
